% =============================================================================
%  lec22.tex —— 第 22 讲 代数（底本 PDF 第 169–174 页）
% =============================================================================
\lecture{22}{代数}
\begin{center}\small\color{black!55}（底本 PDF 第 169–174 页）\end{center}

本讲研究 $\infty$-操作子 $\OO^{\otimes}$ 上的 $\OO$-代数。

\section{代数的 $n$ 元运算}

\nitem\label{L22-1.1} 上次，对任意 $\infty$-操作子 $\OO^{\otimes}$ 与 $\cC^{\otimes}$，我们把 $\cC$ 中的 $\OO$-代数定义为 $\infty$-操作子映射 $p\colon\OO^{\otimes}\to\cC^{\otimes}$，即一个定义在 $\mathrm{Fin}_{*}$ 上、保持惰性态射的函子。我们先来说明这个定义的合理性。

\yizhu{对 $\mathrm{Fin}_{*}$ 中的态射 $\alpha\colon\langle n\rangle\to\langle m\rangle$，若对每个非基点 $j$ 都有 $\alpha^{-1}(\{j\})$ 为单点集，则称 $\alpha$ 为\textbf{惰性}的；若 $\alpha^{-1}(\{*\})=\{*\}$，则称 $\alpha$ 为\textbf{活动}的。$\mathrm{Fin}_{*}$ 中唯一的活动态射 $\langle n\rangle\to\langle1\rangle$ 记作 $\mult_{n}$。\textbf{[来源：B]}}

\nitem\label{L22-1.2} 设 $\cC^{\otimes}\to\mathrm{Fin}_{*}$ 是对称幺半范畴，$p\colon\OO^{\otimes}\to\cC^{\otimes}$ 是 $\cC$ 中的一个 $\OO$-代数。对 $\OO^{\otimes}$ 中任意 $n$ 元运算，即位于 $\mathrm{Fin}_{*}$ 中唯一活动态射 $\mult_{n}\colon\langle n\rangle\to\langle1\rangle$ 上方的活动态射 $f\colon X=(X_{1},\dots,X_{n})\to Y$，我们可以如下构造 $\cC$ 中的一个态射
\begin{equation}\label{L22-eq1.1}
p(X_{1})\otimes\cdots\otimes p(X_{n})\to p(Y).
\end{equation}
由于 $\cC^{\otimes}\to\mathrm{Fin}_{*}$ 是余笛卡尔纤维化，$p(f)$ 有本质唯一的分解 $p(X)\xrightarrow{g}Z\xrightarrow{h}p(Y)$，其中 $g$ 是活动态射 $\mult_{n}$ 上方的余笛卡尔箭头，而 $h$ 落在 $\cC^{\otimes}_{\langle1\rangle}\defeq\cC$ 中。利用 $p$ 保持惰性态射这一假设，可以证明 $p(X)\simeq(p(X_{1}),\dots,p(X_{n}))$。由此
\[
Z\simeq\mult_{n,\dagger}(p(X_{1}),\dots,p(X_{n}))\simeq p(X_{1})\otimes\cdots\otimes p(X_{n}),
\]
其中最后一个同构可经由拉直看出。现在，态射 $h\colon Z\to p(Y)$ 就给出所求的态射 \eqref{L22-eq1.1}。

\begin{exercise}\label{L22-1.3}
把上述构造推广，得到一个函子
\[
\OO^{\otimes}_{\mathrm{active}}\to\cC,\qquad (X_{1},\dots,X_{n})\mapsto p(X_{1})\otimes\cdots\otimes p(X_{n}),
\]
其中 $\OO^{\otimes}_{\mathrm{active}}$ 是 $\OO^{\otimes}$ 的由活动态射构成的 1-全子 $\infty$-范畴。
\end{exercise}

\yizhu{“1-全子 $\infty$-范畴”指这样的子 $\infty$-范畴：它与母范畴有相同的对象，1-态射只取指定的某一类态射，并包含由这些态射张成的全部高阶单形。\textbf{[来源：B]}}

\section{$\mathrm{Triv}$-代数与 $E_{0}$-代数}

\nitem\label{L22-2.1} 上次我们引入了通常的操作子 $\mathrm{Assoc}^{\otimes}$ 与 $\mathrm{Comm}^{\otimes}$。我们将在以后的讲次中详细研究相应的代数与模。眼下先考虑一些更简单的操作子。

\begin{exercise}\label{L22-2.2}
设 $\mathrm{Triv}^{\otimes}$ 是 $\mathrm{Fin}_{*}$ 的由惰性态射构成的子范畴。证明 $\mathrm{Triv}^{\otimes}$ 是一个操作子，并且对任意 $\cC^{\otimes}\in\mathrm{Op}_{\infty}$，在 $\langle1\rangle\in\mathrm{Triv}^{\otimes}$ 处的求值诱导一个等价
\[
\Alg_{\mathrm{Triv}}(\cC)\xrightarrow{\ \simeq\ }\cC.
\]
\end{exercise}

\begin{exercise}\label{L22-2.3}
设 $E_{0}^{\otimes}$ 是 $\mathrm{Fin}_{*}$ 的由满足下述条件的态射 $f\colon\langle n\rangle\to\langle m\rangle$ 构成的子范畴：$f^{-1}\{j\}$ 要么为空、要么是单点集。证明：
\begin{enumerate}[label=(\arabic*),leftmargin=2.4em,itemsep=0.1em]
\item 函子 $E_{0}^{\otimes}\to\mathrm{Fin}_{*}$ 使 $E_{0}^{\otimes}$ 成为一个操作子。
\item 对 $E_{0}^{\otimes}$ 上的任意 $\infty$-操作子 $\cC^{\otimes}$，在 $E_{0}^{\otimes}$ 中唯一态射 $\langle0\rangle\to\langle1\rangle$ 处的求值诱导一个等价
\[
\Alg_{/E_{0}}(\cC)\xrightarrow{\ \simeq\ }\Fun_{E_{0}^{\otimes}}(\Delta^{1},\cC^{\otimes}).
\]
\item 知道一个 $E_{0}$-幺半 $\infty$-范畴 $\cC^{\otimes}$，等价于知道一个 $\infty$-范畴 $\cC$ 连同其中的一个对象 $\mathbf{1}$。更确切地说，有一个由沿态射 $\langle0\rangle\to\langle1\rangle$ 的协变传输给出的等价
\[
\Mon_{E_{0}}(\Cat_{\infty})\xrightarrow{\ \simeq\ }(\Cat_{\infty})_{*}\defeq(\Cat_{\infty})_{\{*\}/}.
\]
\item 对 $E_{0}$-幺半 $\infty$-范畴 $\cC^{\otimes}$，有 $\Alg_{/E_{0}}(\cC)\simeq\cC_{\mathbf{1}/}$。
\end{enumerate}
\end{exercise}

\begin{remark}\label{L22-2.4}
在以后的讲次中，我们将往序列
\[
\mathrm{Triv}^{\otimes}\to E_{0}^{\otimes}\to\mathrm{Assoc}^{\otimes}\to\mathrm{Comm}^{\otimes}
\]
中填入更多的项。
\end{remark}

\section{单位性}

\nitem\label{L22-3.1} 练习~\ref{L22-2.2} 与练习~\ref{L22-2.3} 提示 $E_{0}^{\otimes}$ 是 $\mathrm{Triv}^{\otimes}$ 的带单位版本。二者的差别可以通过下述练习更好地表述。

\begin{exercise}\label{L22-3.2}
证明 $\mathrm{Triv}^{\otimes}$ 中没有零元运算，而 $E_{0}^{\otimes}$ 中有唯一的零元运算。
\end{exercise}

\nitem\label{L22-3.3} 这促使我们作下述定义。

\begin{definition}\label{L22-3.4}
设 $\OO^{\otimes}$ 是 $\infty$-操作子，$\langle0\rangle\in\OO^{\otimes}_{\langle0\rangle}$ 是本质唯一的对象。
\begin{itemize}[leftmargin=2.2em,itemsep=0.1em]
\item 称 $\OO^{\otimes}$ 是\textbf{带单位的}，若对任意颜色 $X\in\OO$，映射空间 $\operatorname{Maps}_{\OO^{\otimes}}(\langle0\rangle,X)$ 可缩。
\item 称 $\OO^{\otimes}$ 是\textbf{非单位的}，若对任意颜色 $X\in\OO$，映射空间 $\operatorname{Maps}_{\OO^{\otimes}}(\langle0\rangle,X)$ 为空。
\end{itemize}
设 $\mathrm{Op}_{\infty}^{\mathrm{unit}}$（相应地，$\mathrm{Op}_{\infty}^{\mathrm{nu}}$）为由带单位（相应地，非单位）的小 $\infty$-操作子构成的全子 $\infty$-范畴。
\end{definition}

\begin{warning}\label{L22-3.5}
$\infty$-操作子可能既不是带单位的，也不是非单位的。
\end{warning}

\begin{warning}\label{L22-3.6}
一个（对称）幺半 $\infty$-范畴 $\cC^{\otimes}$ 作为 $\infty$-操作子未必带单位。$\cC^{\otimes}$ 作为 $\infty$-操作子是否带单位的判别准则见练习~\ref{L22-3.13}。
\end{warning}

\nitem\label{L22-3.7} Lurie 把对象 $\langle0\rangle\in\OO^{\otimes}_{\langle0\rangle}$ 记作 $\varnothing$。但要注意，它一般并不是整个 $\OO^{\otimes}$ 中的始对象。尽管如此，我们有下述结果。

\begin{exercise}\label{L22-3.8}
设 $\OO^{\otimes}$ 是任意 $\infty$-操作子，证明 $\langle0\rangle\in\OO^{\otimes}$ 是终对象。
\end{exercise}

\begin{exercise}\label{L22-3.9}
$\infty$-操作子 $\OO^{\otimes}$ 带单位，当且仅当 $\langle0\rangle\in\OO^{\otimes}$ 是始对象。
\end{exercise}

\begin{definition}\label{L22-3.10}
设 $\OO^{\otimes}$ 是带单位的 $\infty$-操作子，$p\colon\cC^{\otimes}\to\OO^{\otimes}$ 是 $\OO$-幺半 $\infty$-范畴。称 $\cC^{\otimes}$ 中的态射 $Y_{0}\to Y$ 把 $Y$ \textbf{展示}为 $\cC$ 的 \textbf{$X$-单位对象}，若它是位于 $\langle0\rangle\to X$ 上方的 $p$-余笛卡尔边。
\end{definition}

\nitem\label{L22-3.11} 由于 $\langle0\rangle\in\OO^{\otimes}$ 是始对象，且 $\cC^{\otimes}_{\langle0\rangle}\simeq\cC^{\times0}$ 可缩，对固定的 $X\in\cC^{\otimes}$，上述态射 $Y_{0}\to Y$ 本质唯一，且 $Y_{0}\simeq\langle0\rangle$。因此，谈论 $\cC$ 的 $X$-单位对象 $\mathbf{1}_{X}$、并把余笛卡尔边 $\langle0\rangle\to\mathbf{1}_{X}$ 视为隐含的数据，是有意义的。

\begin{warning}\label{L22-3.12}
对象 $\langle0\rangle\in\cC^{\otimes}$ 未必是始对象，余笛卡尔边 $\langle0\rangle\to\mathbf{1}_{X}$ 也并非本质唯一。事实上，这样的边构成的空间可以等同于 $\mathbf{1}_{X}$ 上的自同构构成的空间（练习）。
\end{warning}

\begin{exercise}\label{L22-3.13}
设 $\OO^{\otimes}$ 是带单位的 $\infty$-操作子。证明：$\OO$-幺半 $\infty$-范畴 $\cC^{\otimes}$ 作为 $\infty$-操作子带单位，当且仅当对任意颜色 $X\in\OO$，对象 $\mathbf{1}_{X}\in\cC^{\otimes}_{X}$ 是始对象。
\end{exercise}

\nitem\label{L22-3.14} 在经典范畴论中，对一个（对称）幺半范畴 $\cC$，其中有一个典范的（交换）代数对象，其底对象是幺半单位 $\mathbf{1}$。此外，这个（交换）代数在 $\cC$ 中所有（交换）代数构成的范畴中是始对象。下述结果把它推广到任意 $\infty$-操作子。

\begin{theoremdef}\label{L22-3.15}（HA.3.2.1.8\footnote{Lurie 证明了 (1)$\Leftrightarrow$(2)。等价 (2)$\Leftrightarrow$(3) 由 $\infty$-操作子的公理推出。}）
设 $\OO^{\otimes}$ 是带单位的 $\infty$-操作子，$p\colon\cC^{\otimes}\to\OO^{\otimes}$ 是 $\OO$-幺半 $\infty$-范畴。对一个 $\OO$-代数 $A\in\Alg_{/\OO}(\cC)$，下述条件等价：
\begin{enumerate}[label=(\arabic*),leftmargin=2.4em,itemsep=0.1em]
\item 对象 $A$ 在 $\Alg_{/\OO}(\cC)$ 中是始对象；
\item 对任意对象 $X\in\OO^{\otimes}$，态射 $A(\langle0\rangle)\to A(X)$ 把 $A(X)$ 展示为 $\cC$ 的 $X$-单位对象；
\item 对任意颜色 $X\in\OO$，态射 $A(\langle0\rangle)\to A(X)$ 把 $A(X)$ 展示为 $\cC$ 的 $X$-单位对象。
\end{enumerate}
我们称 $A$ 为 $\cC$ 中的\textbf{单位 $\OO$-代数}\footnote{Lurie 称之为 $\cC$ 中的平凡 $\OO$-代数。}，并记作 $\mathbf{1}_{\cC}$。
\end{theoremdef}

\yizhu{HA 指 Lurie 的 \emph{Higher Algebra}；文中形如 HA.3.2.1.8 的编号即该书中的结果编号。\textbf{[来源：A]}}

\nitem\label{L22-3.16} 我们再把下述结果记录下来。

\begin{proposition}\label{L22-3.17}（HA.2.3.1.9）
嵌入 $\mathrm{Op}_{\infty}^{\mathrm{unit}}\subseteq\mathrm{Op}_{\infty}$ 有左伴随。
\end{proposition}

\begin{remark}\label{L22-3.18}
在以后的讲次中，我们将给 $\mathrm{Op}_{\infty}$ 赋予对称幺半结构，即构造一个 $\infty$-操作子的余笛卡尔纤维化 $\mathrm{Op}_{\infty}^{\otimes}\to\mathrm{Comm}^{\otimes}$，它在 $\langle1\rangle$ 上的纤维是 $\mathrm{Op}_{\infty}$。$\mathrm{Op}_{\infty}$ 的幺半单位由 $\mathrm{Triv}^{\otimes}$ 给出。此外，由 HA.2.3.1.9，复合
\[
\mathrm{Op}_{\infty}\to\mathrm{Op}_{\infty}^{\mathrm{unit}}\to\mathrm{Op}_{\infty}
\]
等同于函子 $E_{0}\otimes-$。于是，$\infty$-操作子 $\OO^{\otimes}$ 带单位，当且仅当
\[
\OO^{\otimes}\to\mathrm{Triv}^{\otimes}\otimes\OO^{\otimes}\to E_{0}^{\otimes}\otimes\OO^{\otimes}
\]
是等价。特别地，$\mathrm{Op}_{\infty}^{\mathrm{unit}}$ 从 $\mathrm{Op}_{\infty}$ 继承对称幺半结构，其幺半单位是 $E_{0}^{\otimes}$。
\end{remark}

\section{非单位代数}

\nitem\label{L22-4.1} 给定 $\infty$-操作子 $\OO^{\otimes}$，我们可以构造一个满足某种泛性质的非单位 $\infty$-操作子 $\OO^{\otimes}_{\mathrm{nu}}$。

\begin{exercise}\label{L22-4.2}
设 $\mathrm{Fin}_{*}^{\mathrm{surj}}\subseteq\mathrm{Fin}_{*}$ 是由满态射构成的子范畴。证明：对任意 $\infty$-操作子 $\OO^{\otimes}$，纤维积
\[
\OO^{\otimes}_{\mathrm{nu}}\defeq\OO^{\otimes}\times_{\mathrm{Fin}_{*}}\mathrm{Fin}_{*}^{\mathrm{surj}}
\]
是非单位 $\infty$-操作子。
\end{exercise}

\begin{exercise}\label{L22-4.3}
证明 $(E_{0}^{\otimes})_{\mathrm{nu}}\simeq\mathrm{Triv}^{\otimes}$。
\end{exercise}

\begin{exercise}\label{L22-4.4}
证明函子
\[
\mathrm{Op}_{\infty}\to\mathrm{Op}_{\infty}^{\mathrm{nu}},\qquad \OO^{\otimes}\mapsto\OO^{\otimes}_{\mathrm{nu}}
\]
是嵌入函子的右伴随。换言之，对任意非单位 $\cD^{\otimes}$，有
\[
\Alg_{\cD}(\OO_{\mathrm{nu}})\xrightarrow{\ \simeq\ }\Alg_{\cD}(\OO).
\]
\end{exercise}

\begin{definition}\label{L22-4.5}
设
\begin{center}
\begin{tikzcd}[column sep=5.5em,row sep=2.8em]
{\OO'}^{\otimes}\arrow[dr] & & \cC^{\otimes}\arrow[dl]\\
& \OO^{\otimes} &
\end{tikzcd}
\end{center}
是 $\infty$-操作子之间的映射。$\cC$ 中的一个\textbf{非单位 $\OO'$-代数}（相对于 $\OO$）是
\[
\Alg^{\mathrm{nu}}_{{\OO'}/\OO}(\cC)\defeq\Alg_{{\OO'}_{\mathrm{nu}}/\OO}(\cC)
\]
中的一个对象。
\end{definition}

\nitem\label{L22-4.6} 回忆：对一个非单位的（交换）环，是否带单位是一条性质而非结构：单位元是唯一的。这一点可以推广到 $\infty$-情形：

\begin{theorem}\label{L22-4.7}（HA.5.4.4.5, 5.4.4.7）
设 $\OO^{\otimes}$ 是 $\mathrm{Assoc}^{\otimes}$ 或 $\mathrm{Comm}^{\otimes}$，$\cC^{\otimes}\to\OO^{\otimes}$ 是 $\OO$-幺半 $\infty$-范畴。则 $\OO^{\otimes}_{\mathrm{nu}}\to\OO^{\otimes}$ 诱导一个全忠实函子
\[
\Alg_{/\OO}(\cC)\to\Alg_{/\OO_{\mathrm{nu}}}(\cC).
\]
\end{theorem}

\begin{remark}\label{L22-4.8}
事实上，前引文献证明了：上述嵌入的本质像恰好由 $\cC$ 中那些在 $\mathrm{h}\cC$ 中的像带单位的非单位结合（相应地，交换）代数组成。
\end{remark}

\begin{remark}\label{L22-4.9}
对 $1\le k\le\infty$ 的 $\infty$-操作子 $E_{k}^{\otimes}$，类似的论断仍然成立。
\end{remark}

\begin{warning}\label{L22-4.10}
对一般的 $\infty$-操作子，类似的论断不成立。
\end{warning}

\begin{exercise}\label{L22-4.11}
证明：对一个 $E_{0}$-幺半 $\infty$-范畴 $\cC$，$\Alg_{/E_{0}}(\cC)\to\Alg_{/(E_{0})_{\mathrm{nu}}}(\cC)$ 未必全忠实。
\end{exercise}

\section{代数的极限与余极限}

\nitem\label{L22-5.1} 本节研究 $\OO$-幺半 $\infty$-范畴 $\cC^{\otimes}\to\OO^{\otimes}$ 中 $\OO$-代数的极限与余极限。

\nitem\label{L22-5.2} 在经典代数中，环的图表 $i\mapsto A_{i}$ 的底层 Abel 群的极限 $\lim A_{i}$ 带有环结构，该结构由
\[
\Bigl(\lim_{i\in K}A_{i}\Bigr)\otimes\Bigl(\lim_{i\in K}A_{i}\Bigr)\to\lim_{(i,j)\in K^{2}}(A_{i}\otimes A_{j})\to\lim_{i\in K}(A_{i}\otimes A_{i})\to\lim_{i\in K}A_{i}
\]
给出。这可以推广到任意 $\infty$-操作子 $\OO^{\otimes}$。

\begin{theorem}\label{L22-5.3}（HA.3.2.2.5）
设 $\cC^{\otimes}\to\OO^{\otimes}$ 是 $\OO$-幺半 $\infty$-范畴，使得对某个固定的指标单纯集 $K$，每个纤维 $\cC^{\otimes}_{X}$ 都有 $K$-指标极限。则 $\infty$-范畴 $\Alg_{/\OO}(\cC)$ 有 $K$-指标极限，且遗忘函子
\[
\Alg_{/\OO}(\cC)\to\cC^{\otimes}_{X},\qquad X\in\OO
\]
保持并检测 $K$-指标极限。
\end{theorem}

\begin{remark}\label{L22-5.4}
注意，在上述定理中，我们并不假设协变传输函子 $\cC^{\otimes}_{X}\to\cC^{\otimes}_{Y}$ 保持 $K$-指标极限。
\end{remark}

\nitem\label{L22-5.5} 正如我们在 $K=\varnothing$ 的情形所见，定理~\ref{L22-5.3} 关于余极限的类似论断一般不成立。事实上，上述复合会变成一个对应
\[
\Bigl(\colim_{i\in K}A_{i}\Bigr)\otimes\Bigl(\colim_{i\in K}A_{i}\Bigr)\leftarrow\colim_{(i,j)\in K^{2}}(A_{i}\otimes A_{j})\leftarrow\colim_{i\in K}(A_{i}\otimes A_{i})\to\colim_{i\in K}A_{i}.
\]
这提示：若两个向左的态射可逆，类似的论断就会成立。为给出精确的陈述，我们需要一些定义。

\begin{definition}\label{L22-5.6}
设 $\OO^{\otimes}$ 是任意 $\infty$-操作子，$p\colon\cC^{\otimes}\to\OO^{\otimes}$ 是 $\OO$-幺半范畴。对单纯集 $K$，称 $p$ \textbf{与 $K$-指标（余）极限相容}，若
\begin{itemize}[leftmargin=2.2em,itemsep=0.1em]
\item 对任意对象 $X\in\OO^{\otimes}$，纤维 $\cC^{\otimes}_{X}$ 有 $K$-指标（余）极限；
\item 对 $\OO^{\otimes}$ 中任意态射 $X\to X'$，协变传输函子 $\cC^{\otimes}_{X}\to\cC^{\otimes}_{X'}$ 保持 $K$-指标（余）极限。
\end{itemize}
称 $\cC$ 是\textbf{可表现 $\OO$-幺半 $\infty$-范畴}，若每个 $\cC^{\otimes}_{X}$ 可表现，且 $p$ 与小余极限相容。
\end{definition}

\begin{theorem}\label{L22-5.7}（HA.3.2.3.1）
设 $K$ 是筛单纯集，$\cC^{\otimes}\to\OO^{\otimes}$ 是与 $K$-指标余极限相容的 $\OO$-幺半 $\infty$-范畴。则 $\infty$-范畴 $\Alg_{/\OO}(\cC)$ 有 $K$-指标余极限，且遗忘函子
\[
\Alg_{/\OO}(\cC)\to\cC^{\otimes}_{X},\qquad X\in\OO
\]
保持并检测 $K$-指标余极限。
\end{theorem}

\nitem\label{L22-5.8} 对一般的余极限，第一个论断仍然成立，但第二个论断将会失效。

\begin{theorem}\label{L22-5.9}（HA.3.1.3.5, 3.2.3.3）
设 $\cC^{\otimes}\to\OO^{\otimes}$ 是与小余极限相容的 $\OO$-幺半 $\infty$-范畴，并设 $\OO^{\otimes}$ 本质小。则
\begin{enumerate}[label=(\roman*),leftmargin=2.6em,itemsep=0.1em]
\item 遗忘函子
\[
\Alg_{/\OO}(\cC)\to\Fun_{\OO}(\OO,\cC)
\]
有左伴随
\[
\mathrm{Free}_{\OO}\colon\Fun_{\OO}(\OO,\cC)\to\Alg_{/\OO}(\cC).
\]
\item $\infty$-范畴 $\Alg_{/\OO}(\cC)$ 有小余极限。
\end{enumerate}
\end{theorem}

\begin{remark}\label{L22-5.10}
如同经典代数中那样，$\OO$-代数的有限余积可以通过用自由 $\OO$-代数对它们作单纯消解来计算。确切地说，对一个 $\OO$-代数 $A$，有单纯对象
\[
\Delta^{\mathrm{op}}\to\Alg_{/\OO}(\cC),\qquad [n]\mapsto\mathrm{Free}_{\OO}^{\circ(n+1)}(A),
\]
其余极限等价于 $A$。这一点将在以后的讲次中讨论。
\end{remark}

\begin{remark}\label{L22-5.11}
在以后的讲次中，我们将把 $\mathrm{Comm}$-代数的有限余积等同于它们的张量积。
\end{remark}

\nitem\label{L22-5.12} 我们再把下述结果记录下来。

\begin{proposition}\label{L22-5.13}（HA.3.2.3.5）
设 $\cC^{\otimes}\to\OO^{\otimes}$ 是可表现 $\OO$-幺半 $\infty$-范畴，且 $\OO^{\otimes}$ 本质小。则 $\Alg_{/\OO}(\cC)$ 可表现。
\end{proposition}

\appendixsection{增广代数}

\nitem\label{L22-A.1} 回忆：对一个非单位的（交换）环 $A$，$\Z\oplus A$ 上有一个典范的带单位（交换）环结构。此外，这个构造给出非单位（交换）环及其模的理论与增广（交换）环及其模的理论之间的一个等价。这些可以推广到任意带单位的 $\infty$-操作子 $\OO^{\otimes}$ 以及满足某些条件的 $\OO$-幺半 $\infty$-范畴 $\cC^{\otimes}$。

\nitem\label{L22-A.2} \textbf{建议阅读}\quad HA.5.4.4。
